\documentclass[conference]{IEEEtran}
\IEEEoverridecommandlockouts
\usepackage{cite}
\usepackage{amsmath,amssymb,amsfonts}
\usepackage{algorithmic}
\usepackage{graphicx}
\usepackage{textcomp}
\usepackage{xcolor}
\def\BibTeX{{\rm B\kern-.05em{\sc i\kern-.025em b}\kern-.08em
    T\kern-.1667em\lower.7ex\hbox{E}\kern-.125emX}}
\begin{document}

\title{Structured Intelligence from Unstructured Web Listings: A Multi-Source Data Engineering Pipeline for Indian Wedding Venue Discovery, Normalization, and Market Analysis\\
{\footnotesize \textsuperscript{*}Note: Sub-titles are not captured in Xplore and
should not be used}
\thanks{Prepared as part of the 6th Semester Engineering Project, 2026.}
}

\author{\IEEEauthorblockN{1\textsuperscript{st} Given Name Surname}
\IEEEauthorblockA{\textit{dept. name of organization (of Aff.)} \\
\textit{name of organization (of Aff.)}\\
City, Country \\
email address or ORCID}
\and
\IEEEauthorblockN{2\textsuperscript{nd} Given Name Surname}
\IEEEauthorblockA{\textit{dept. name of organization (of Aff.)} \\
\textit{name of organization (of Aff.)}\\
City, Country \\
email address or ORCID}
\and
\IEEEauthorblockN{3\textsuperscript{rd} Given Name Surname}
\IEEEauthorblockA{\textit{dept. name of organization (of Aff.)} \\
\textit{name of organization (of Aff.)}\\
City, Country \\
email address or ORCID}
\and
\IEEEauthorblockN{4\textsuperscript{th} Given Name Surname}
\IEEEauthorblockA{\textit{dept. name of organization (of Aff.)} \\
\textit{name of organization (of Aff.)}\\
City, Country \\
email address or ORCID}
\and
\IEEEauthorblockN{5\textsuperscript{th} Given Name Surname}
\IEEEauthorblockA{\textit{dept. name of organization (of Aff.)} \\
\textit{name of organization (of Aff.)}\\
City, Country \\
email address or ORCID}
\and
\IEEEauthorblockN{6\textsuperscript{th} Given Name Surname}
\IEEEauthorblockA{\textit{dept. name of organization (of Aff.)} \\
\textit{name of organization (of Aff.)}\\
City, Country \\
email address or ORCID}
}

\maketitle

\begin{abstract}
India's wedding market sits at roughly 3.75 lakh crore rupees -- about USD 45 billion according to the KPMG 2023 report \cite{b1}. Prime Minister Modi, during the 107th episode of Mann Ki Baat in November 2023, urged families to stop going abroad for weddings and instead celebrate on Indian soil so that ``the country's money stays in the country'' \cite{b22}. The Government followed up with the ``Wed in India'' campaign along with FICCI, positioning India as a global wedding destination. But here is the disconnect -- despite all the policy push, there is no publicly available structured dataset of Indian wedding venues. Not one. Platforms like WedMeGood, VenueLook, ShaadiSaga hold all the data behind proprietary walls; prices show up in six or more text formats ranging from per-plate rates to Lakh ranges. Cross-platform comparison without automated tools is basically impossible.

This paper presents WedVenueIQ, a five-phase pipeline that scrapes venue listings from three major sources, normalizes messy Indian price strings through a six-pass regex cascade, deduplicates cross-source records using city-gated fuzzy matching, validates everything against a schema, then generates market intelligence on top. The Rajasthan pilot produced 2,036 unique cleaned venues from 5,942 raw records. Price coverage hit 99.9 percent post-imputation. We also introduce three original frameworks -- the Wedding Venue Value Index, Cost-Popularity Z-score Classification, and Micro-Location Composite Score -- plus a specialized analysis of the NRI overseas wedding segment worth an estimated USD 5 billion.
\end{abstract}

\begin{IEEEkeywords}
data engineering, web scraping, information extraction, price normalization, record linkage, fuzzy deduplication, Indian wedding industry, Wed in India, market analytics, hospitality data
\end{IEEEkeywords}

\section{Introduction}

\subsection{The Wedding Economy and National Policy Context}

India sees around 10 million weddings a year. Of those, an estimated 2.5 million qualify as destination events -- held at heritage palaces, lakeside resorts, beach properties, farmhouses \cite{b1}. The KPMG 2023 report pegs the overall market at 3.75 lakh crore INR with destination weddings growing at 25 percent CAGR through 2028. FICCI, in collaboration with the Ministry of Tourism, has been pushing wedding tourism as a pillar of India's hospitality strategy \cite{b23}.

The political push started getting serious in late 2023. On November 26, during Mann Ki Baat Episode 107, PM Modi said -- ``these days a new milieu is being created by some families to go abroad and conduct weddings. Is this at all necessary? If we celebrate the festivities of marriages on Indian soil, the country's money will remain in the country'' \cite{b22}. He later framed it as ``Wed in India, like Made in India, Marry in India'' at a Gujarat event in January 2024. The Ministry of Tourism responded with the ``India Says I Do'' campaign -- a 360-degree digital and offline initiative profiling about 25 key destinations across states including Rajasthan, Goa, Kerala, and Uttarakhand \cite{b24}. FICCI co-organized the ``Wed in India'' expo alongside the Great India Travel Bazaar (GITB) in Jaipur, May 2024, bringing together planners, state tourism boards, and operators.

But the gap between policy ambition and data infrastructure is wide. None of these campaigns are backed by a structured, queryable database of venues. The tourism.gov.in website lists destinations but provides no machine-readable venue data -- no prices, no capacities, no ratings. State tourism boards promote their heritage properties through brochures and Instagram posts but offer nothing a researcher or planner can query programmatically. The data simply does not exist in structured form.

\subsection{Weddings as Soft Power}

There is a broader angle here worth noting. Destination weddings are increasingly recognized as instruments of cultural soft power -- when you sell your culture through celebrations, you reshape how the world perceives India \cite{b21}. The optics of a royal wedding at a Rajasthani palace or a backwater ceremony in Kerala circulate globally through social media. This is not just tourism; it is cultural diplomacy. Yet the infrastructure to measure, benchmark, or optimize this soft power asset remains non-existent from a data perspective.

\subsection{What Has Been Done and What Is Lagging}

The Government of India and state tourism departments have taken concrete steps. Rajasthan aggressively markets its palace circuit. Kerala positions itself for eco-luxury weddings. Goa leverages its beach infrastructure. Uttarakhand promotes Himalayan backdrops. The Ministry runs digital campaigns, collaborates with influencers, participates in global travel expos.

What is lagging: the digital backbone. The tourism.gov.in portal has no venue search functionality beyond basic destination pages. No filtering by price, capacity, or type. No comparison tools. No structured data exports. Instagram campaigns generate engagement but zero queryable intelligence. A wedding planner trying to find all palace venues in Rajasthan under 15 lakhs with capacity above 300 guests has to manually visit three or four separate websites, copy data into spreadsheets, and eyeball duplicates. That is the state of things in 2024.

Our suggestion -- and what this paper demonstrates -- is that a data engineering pipeline can bridge this gap. WedVenueIQ is not a replacement for government campaigns; it is the missing data layer underneath them.

\begin{table}[htbp]
\caption{State-Wise ``Wed in India'' Campaign Activities}
\begin{center}
\begin{tabular}{|l|l|l|}
\hline
\textbf{State} & \textbf{Focus} & \textbf{Promotion Channel} \\
\hline
Rajasthan & Palace circuit, forts & GITB Expo, Insta, YT \\
\hline
Goa & Beach resorts & Digital ads, influencers \\
\hline
Kerala & Backwaters, eco-luxury & State tourism portal \\
\hline
Uttarakhand & Himalayan backdrops & PM quote, digital push \\
\hline
Maharashtra & Urban luxury venues & Event expos \\
\hline
Karnataka & Heritage, hill stations & State board brochures \\
\hline
\end{tabular}
\label{tab:statewise}
\end{center}
\end{table}

\begin{table}[htbp]
\caption{``India Says I Do'' Campaign -- Digital Presence Audit (2024)}
\begin{center}
\begin{tabular}{|l|l|l|}
\hline
\textbf{Platform} & \textbf{Content Type} & \textbf{Status} \\
\hline
YouTube (MoT) & Destination reels & Active, low views \\
\hline
Instagram (Incredible India) & Wedding destination posts & Active, organic reach \\
\hline
tourism.gov.in & Static destination pages & No venue search/filter \\
\hline
PIB Press Releases & Expo coverage, policy & Regular updates \\
\hline
State Tourism Portals & Brochure PDFs, galleries & No structured data \\
\hline
FICCI/GITB Expos & Offline B2B matchmaking & Annual, Jaipur-based \\
\hline
\end{tabular}
\label{tab:campaign}
\end{center}
\end{table}


\subsection{Problem Statement}

Despite the economic scale, the venue discovery landscape is fragmented badly. WedMeGood, VenueLook, ShaadiSaga operate as closed silos. Data quality varies wildly across listings. Pricing for the same venue might show as a ``starting price'' on one platform and a ``per-plate rate'' on another. Without normalization, cross-platform comparison is computationally impossible for any consumer. This creates measurable inefficiency -- venues in low-competition areas keep prices high because there is no benchmark; couples overspend without knowing about comparable alternatives in smaller cities.

\subsection{Contributions}

This work contributes the following. A validated dataset of 2,036 unique Indian wedding venues from the Rajasthan pilot -- the first publicly constructed dataset of its kind. A six-pass normalization algorithm for Indian currency units achieving 94 percent coverage. A city-gated deduplication engine removing 16.3 percent cross-source duplicates. Three original market intelligence frameworks. Plus an analysis of the NRI overseas wedding market segment. The pipeline is modular and configuration-driven; extending it to a new state requires editing a JSON file, not code.

\section{Literature Review}

\subsection{Web Scraping and Data Extraction}

Pulling structured data from web pages through CSS selectors and XPath is well-documented \cite{b3}. But modern sites run on JavaScript SPAs -- they render content client-side, which means static HTTP parsing returns empty shells. Headless browsers like Selenium \cite{b4} and Playwright \cite{b5} solve this but introduce latency and complexity. Anti-bot countermeasures including rate limiting, IP fingerprinting, JS-based detection get handled via User-Agent rotation, random delays, exponential backoff on HTTP 429 \cite{b6}.

Schema.org JSON-LD was supposed to make things easier by embedding machine-readable metadata in HTML \cite{b7}. In practice, our audit of Indian wedding platforms found fewer than 30 percent of listings had complete Schema.org markup \cite{b8}. So multi-strategy extraction -- JSON-LD where available, HTML fallback otherwise -- is unavoidable.

\subsection{Record Linkage}

Identifying duplicate records across heterogeneous sources is a foundational data integration problem \cite{b9}. The Fellegi-Sunter model \cite{b10} gives the probabilistic basis. String similarity via Levenshtein distance \cite{b11} and token-set variants \cite{b12} are standard practice. Blocking strategies that restrict comparisons to records sharing a key attribute (like city) reduce complexity from $O(N^2)$ to something manageable \cite{b13}. ML-based entity matching \cite{b14} and deep approaches like DeepMatcher \cite{b15} exist but need labeled training data. We do not have that here, so threshold-based fuzzy matching was the right call.

\subsection{Price Normalization}

Extracting numbers from free-text price strings is an active problem \cite{b16}. Rule-based regex works when the vocabulary is finite. Indian pricing is a special case -- Lakh ($10^5$), Crore ($10^7$), per-plate conventions, range expressions crossing units, multiple rupee symbols. No published library handles all of these. Standard tools like Python's price-parser fail on ``2.5L -- 10L'' or ``Rs 800/plate.'' We found nothing in the literature addressing Indian wedding price normalization specifically.

\subsection{Hospitality Analytics}

Hotel analytics is mature -- STR Global \cite{b17} enables ADR and RevPAR benchmarking worldwide. Event venue analytics is far behind. Andersson and Getz \cite{b18} discuss capacity utilization but not pricing. DBSCAN-based spatial clustering has been used in tourism research \cite{b19} but never for Indian wedding venues. Luxury corridor detection through price-density mapping exists for hotels \cite{b20} -- nothing equivalent for weddings.

\subsection{Indian Wedding Research}

Published data on the Indian wedding industry is almost entirely aggregate-level -- KPMG numbers \cite{b1}, BCG estimates. Venue-level data does not exist publicly. Academic work on Indian weddings is overwhelmingly sociological and cultural \cite{b21}. No prior data engineering, information extraction, or market analytics paper exists for this sector. That gap is what this work addresses.

\section{Research Gaps}

\subsection{Gap 1: No Structured Dataset Exists}

Zero. We checked Kaggle, Google Dataset Search, data.gov.in, Zenodo, Harvard Dataverse, IEEE DataPort as of October 2024. Nothing.

\subsection{Gap 2: No Normalization Algorithm for Indian Venue Prices}

Libraries like price-parser are built for Western e-commerce. They choke on Lakhs, Crores, range expressions, per-plate formats. Our pilot found six distinct price format families. No existing tool handles all six.

\subsection{Gap 3: Unquantified Cross-Source Duplication}

Each platform maintains venues independently. Same venue, different spelling, different price, different review count. Nobody has measured the overlap rate. We found 16.3 percent -- enough to badly distort any naive count-based analysis.

\subsection{Gap 4: No Analytical Layer}

No luxury corridor detection. No cost-popularity classification. No micro-location scoring. Cannot build these without Gap 1 being solved first.

\subsection{Gap 5: No Reproducible Collection Method}

Ad-hoc scraping scripts are brittle, undocumented, source-coupled. No published methodology for systematic multi-source Indian hospitality data collection. Our pipeline uses a cities.json config that decouples geography from source logic.


\section{System Architecture}

WedVenueIQ runs as a five-phase linear pipeline. Each phase is its own Python module with defined input and output schemas -- you can swap or rewrite any phase without breaking the others.

The phases: (1) multi-source scraping, (2) cleaning and normalization, (3) schema validation plus imputation, (4) structured export, (5) analytical intelligence. Phases 1 through 4 are the data engineering core. Phase 5 is the application layer sitting on top.

Geographic config lives in a cities.json file outside the code. Want to add a new state? Edit the JSON. No code changes needed. Currently supports 25 cities across 10 Indian states.

\section{Methodology}

\subsection{Data Collection}

Three Tier-1 sources were picked based on coverage of Rajasthan venues, field richness, and accessibility without login.

WedMeGood -- largest platform by listing count. Uses server-side rendered React with a \texttt{window.\_\_INITIAL\_STATE\_\_} JSON blob embedded in page source. We extract name, pricing, capacity, rating, address directly from that object. No HTML parsing needed.

VenueLook -- second-largest. Has Schema.org JSON-LD tags for name, rating, address. Capacity and pricing sit in HTML comparison tables; extracted via CSS selector alignment.

ShaadiSaga -- React SPA, JavaScript-rendered. Requires headless Chrome via Selenium. Fields come from DOM class selectors after a wait for page load.

All sources share one protocol: check SQLite cache first (7-day TTL), then HTTP GET or browser load, parse, extract to memory, apply random delay (1.5--3.0s), rotate User-Agent every 50 requests, exponential backoff on 429 errors (30s initial, doubles up to 120s).

\begin{table}[htbp]
\caption{Field Extraction Methods by Source}
\begin{center}
\begin{tabular}{|l|c|c|c|}
\hline
\textbf{Field} & \textbf{WedMeGood} & \textbf{VenueLook} & \textbf{ShaadiSaga} \\
\hline
Name & JSON state & JSON-LD & DOM class \\
\hline
Price & JSON state & HTML table & DOM class \\
\hline
Capacity & JSON state & HTML table & DOM class \\
\hline
Rating & JSON state & JSON-LD & DOM class \\
\hline
Reviews & JSON state & HTML span & DOM class \\
\hline
Address & JSON state & JSON-LD & DOM class \\
\hline
\end{tabular}
\label{tab:extraction}
\end{center}
\end{table}

\subsection{Six-Pass Price Normalization}

Raw prices come in at least six format families. The engine runs each string through a regex cascade -- first match wins, processing stops.

Pass 1 catches per-plate patterns -- ``800/plate,'' ``Rs 1,200 per plate,'' ``900 pp.'' Extracted value goes to price\_per\_plate; a 300-guest assumption generates a floor estimate for comparison.

Pass 2 handles Lakh ranges. ``2.5 Lakhs -- 10 Lakhs'' or ``2L to 10L'' -- both endpoints extracted, multiplied by $10^5$.

Pass 3 picks up single Lakh values with all spelling variants -- Lakh, Lakhs, Lac, Lacs. Multiply by $10^5$.

Pass 4 is for Crore values. Palace listings sometimes hit ``1.5 Cr'' territory. Multiply by $10^7$.

Pass 5 handles K-denomination. ``50K -- 80K'' -- multiply both by $10^3$.

Pass 6 is the fallback. Raw numbers, with or without commas. If nothing above matched -- this catches whatever is left.

Validation rejects records where min\_price is zero or negative, max\_price falls below min\_price, or price\_per\_plate exceeds 50,000.

\begin{table}[htbp]
\caption{Normalization Examples: Raw Input to Canonical INR}
\begin{center}
\begin{tabular}{|l|l|c|}
\hline
\textbf{Raw String} & \textbf{Pass} & \textbf{Output (INR)} \\
\hline
``Rs. 850/plate'' & Per-Plate & 2,55,000* \\
\hline
``2.5 Lakh -- 8 Lakh'' & Range & 2,50,000 (min) \\
\hline
``1.2 Crore'' & Denomination & 1,20,00,000 \\
\hline
``45K -- 60K'' & K-Notation & 45,000 (min) \\
\hline
``15,00,000'' & Raw & 15,00,000 \\
\hline
\end{tabular}
\label{tab:norm}
\end{center}
\end{table}

\subsection{Venue Type Classification}

Two-pass classifier. Six categories: palace, resort, beach, farmhouse, hotel, banquet hall.

Pass 1 checks the first token of the raw type label against a keyword dictionary. Quick match returns immediately. Pass 2 concatenates name plus type label, lowercases it, scans for keywords in priority order. Indian heritage terms -- mahal, haveli, garh, kothi -- map to palace. Chain names like Marriott, Oberoi, Taj, Leela map to hotel. First hit wins.

Records matching neither pass get flagged for manual review. About 16.8 percent fell into this bucket in the pilot.

\subsection{City-Gated Deduplication}

Two stages. Stage 1 is within-source: identical URLs removed, identical (name, city) pairs collapsed. Stage 2 is cross-source fuzzy matching with city-gated blocking.

Records grouped by city. Within each group, incoming records compared against all accepted records using token\_set\_ratio from fuzzywuzzy. This algorithm sorts tokens before comparing -- critical for Indian naming conventions where ``Taj Lake Palace'' and ``Lake Palace Taj Resort'' should match. Threshold: 85. Primary record kept is the one with higher rating; null fields filled from the secondary before discarding.

City-gating cuts pairwise comparisons from $N^2 / 2$ globally to $\sum_i n_i^2 / 2$ per city. For 2,036 records that is roughly 13x reduction. No recall loss -- venues in different cities cannot be duplicates.

The 85 threshold was picked empirically from 200 labeled pairs. ``Taj Lake Palace'' vs ``Lake Palace by Taj Hotels'' scores 92 (correct match). ``Taj Hotel Jaipur'' vs ``Hotel Taj Mahal Palace'' scores 78 (correctly kept separate).

\subsection{Schema Validation and Imputation}

Hard rules: name non-empty, city matches cities.json, capacity positive, min\_price non-negative, max\_price at least equal to min\_price, rating between 0.0 and 5.0. Fail any hard rule -- record dropped with reason code logged.

Soft rules flag but do not drop: review count below 5, price above 1 crore, rating missing with 50+ reviews. For missing prices after normalization -- city-level median imputation applied. Imputed records flagged with a boolean for transparency.

\section{Analytical Frameworks}

\subsection{Wedding Venue Value Index}

WVVI ranks venues by consumer utility -- quality, social proof, affordability combined. For venue $v$:

\begin{equation}
WVVI(v) = 0.4 \cdot \hat{r}_v + 0.3 \cdot \hat{c}_v + 0.3 \cdot (1 - \hat{p}_v) \label{eq:wvvi}
\end{equation}

$\hat{r}_v$, $\hat{c}_v$, $\hat{p}_v$ are min-max normalized rating, review count, price. Rating gets 0.4 weight because it is the strongest quality signal. Price is inverted -- lower price means higher value score.

\subsection{Z-Score Classification}

Cities classified as Overpriced, Undervalued, or Balanced using two-dimensional Z-scores:

\begin{equation}
z_p(c) = \frac{\tilde{p}(c) - \mu_p}{\sigma_p}, \quad z_r(c) = \frac{\tilde{r}(c) - \mu_r}{\sigma_r} \label{eq:zscore}
\end{equation}

Overpriced if $z_p > 0.5$ and $z_r < -0.5$. Undervalued if $z_p < -0.5$ and $z_r > 0.5$. Everything else is Balanced.

\subsection{Micro-Location Composite Score}

MLS ranks neighborhoods within a city by multi-factor desirability:

\begin{equation}
MLS(a) = 0.4 \cdot \hat{d}_a + 0.3 \cdot \hat{q}_a + 0.3 \cdot \hat{v}_a \label{eq:mls}
\end{equation}

Density (venue count), average rating, price diversity (range width as budget variety proxy). Density at 0.4 because having options nearby matters practically.

\section{Implementation}

\subsection{Technology Stack}

The full pipeline is built on Python 3.10. For HTTP, we use the requests library (v2.31+) with connection pooling across sessions. HTML gets parsed by BeautifulSoup4 backed by the lxml engine -- faster than the default html.parser by a noticeable margin on large pages. JavaScript-heavy sites require Selenium 4.18+ driving a headless Chrome instance; ChromeDriver versioning is handled automatically by webdriver-manager so the pipeline does not break when Chrome updates.

Data wrangling happens in pandas 2.2+ with numpy underneath. Fuzzy matching runs through fuzzywuzzy, but we compile it with the python-Levenshtein C extension -- without that, token\_set\_ratio is painfully slow on thousands of comparisons. Visualization splits between matplotlib plus seaborn for static publication figures and plotly for the interactive dashboard that ships as a standalone HTML file. Caching uses Python's built-in sqlite3 module. Excel export via openpyxl. Testing with pytest and pytest-cov for coverage tracking.

\subsection{Key Design Decisions}

SQLite was chosen over Redis for the scraping cache. Redis needs a running server process and cannot be distributed as a single portable file. SQLite has zero dependencies, stores everything in one .db file, and handles our write pattern -- batch inserts once a week, reads on every subsequent run -- without issue. A 7-day TTL is implemented as a simple timestamp check on cached responses.

City-gated blocking was chosen over global all-pairs comparison. For 2,036 records, naive all-pairs would mean roughly 2.07 million fuzzy string comparisons. City-gated blocking drops that to about 150,000. A 13x reduction. And the cost is zero -- venues in Jaipur and Udaipur cannot possibly be duplicates of each other, so restricting comparisons within cities loses nothing.

We tested GPT-4 prompt-based price extraction against our rule-based cascade on 100 held-out price strings. The cascade hit 94 percent accuracy at roughly 10 milliseconds per record, no API dependency, no cost. GPT-4 got 89 percent on the same set, at 800ms per record and about USD 0.002 each. For 2,036 records the rule-based approach is faster by two orders of magnitude and free. The choice is obvious.

The scraper architecture uses inheritance. All three source scrapers extend a common BaseScraper class that provides caching, delay logic, and backoff handling. Adding a fourth source -- say Sulekha or Weddingz.in -- means writing one new file with a single scrape\_city() method. Everything else is inherited. This was tested by adding a mock source in under 15 minutes during development.

\section{Results}

\subsection{Pilot Statistics}

Rajasthan pilot covered five cities -- Jaipur, Udaipur, Jodhpur, Pushkar, Jaisalmer. Three sources. 5,942 raw records collected. Within-source dedup brought it to 2,433. Cross-source fuzzy dedup produced 2,036 unique venues. 397 duplicates removed -- 16.3 percent rate. Total runtime about 10 minutes from cold cache.

\begin{table}[htbp]
\caption{Pilot Volume Metrics}
\begin{center}
\begin{tabular}{|l|c|}
\hline
\textbf{Stage} & \textbf{Count} \\
\hline
Raw records & 5,942 \\
\hline
After within-source dedup & 2,433 \\
\hline
After cross-source dedup & 2,036 \\
\hline
Duplicates removed & 397 (16.3\%) \\
\hline
Runtime & ~10 min \\
\hline
\end{tabular}
\label{tab:volume}
\end{center}
\end{table}

\subsection{Field Coverage}

Name and city hit 100 percent -- every record has both. Rating coverage sits at 94.4 percent; we deliberately chose not to impute missing ratings because median-filling a quality signal would inject noise rather than information. Price coverage was 83.8 percent after the six-pass normalization; city-level median imputation pushed that to 99.9 percent. The remaining 0.1 percent are edge cases where even the city median was unavailable (very small cities with fewer than 3 priced listings). Venue type coverage reached 83.2 percent through the two-pass keyword classifier. Area or micro-location data was available for only 62 percent of records -- many listings simply do not include neighborhood-level location.

\subsection{Geographic and Type Distribution}

Jaipur dominates the dataset at 1,085 of 2,036 records -- 53.3 percent. Not surprising; it is the state capital with the densest commercial hospitality infrastructure. Udaipur follows at 431, which reflects its position as the go-to luxury destination wedding city in Rajasthan. Jodhpur contributes 290 records. Pushkar at 121 and Jaisalmer at 92 round out the top five.

By venue type: banquet halls account for 36 percent, the bread-and-butter of Indian wedding hosting. Hotels at 24.8 percent. Resorts at 20.5. Palace venues are only 12.4 percent by count but punch far above their weight in price tier -- the average palace minimum price is more than 3x the average banquet hall price. Farmhouses at 6.1 percent. Beach properties basically do not exist in Rajasthan, at 0.2 percent.

\subsection{Price Intelligence}

The cheapest listing in the dataset was Rs. 175. Almost certainly a data entry error on the source platform -- our soft rules flagged it automatically. The most expensive was Rs. 95 Lakhs for a heritage palace in the City Palace area of Udaipur. Median across all 2,036 venues: Rs. 8.5 Lakhs. That number gives you a reasonable baseline for what a destination wedding venue costs in Rajasthan.

Per-plate pricing ranged from Rs. 400 to Rs. 4,500. The distribution was right-skewed as expected -- a small number of high-end venues pull the mean well above the median.

\subsection{Analytical Results}

The luxury corridor detection algorithm identified Badi Lake and City Palace micro-areas in Udaipur as the densest concentration of premium venues anywhere in the dataset. Seven or more venues with average minimum price above Rs. 10 Lakhs and average rating above 4.6 out of 5.0 sit within a 3-kilometer radius there. That is not a coincidence -- it is a genuine luxury cluster.

Z-score classification flagged Pushkar and Kumbhalgarh as Undervalued. Both show below-median price Z-scores combined with above-median review engagement Z-scores. Meaning: strong customer interest at prices well below the Rajasthan average. For couples on a budget who still want a destination wedding, these cities offer real value. And for tourism boards -- Pushkar in particular deserves more promotion spend than it currently receives.

MLS ranked City Palace area in Udaipur first across all micro-locations. It combines the highest venue density, above-average quality scores, and the widest price range diversity -- meaning couples at every budget level can find something there. But the metric also surfaced micro-locations in Jaipur (Amer area, Mansarovar) and Jodhpur (Umaid zone) as serious alternatives that nobody talks about. That is the value of quantitative ranking over brand-driven perception.

\subsection{State-Wise Campaign Alignment}

Our structured data maps directly onto the Government's ``Wed in India'' state promotions and exposes alignment gaps. Rajasthan's palace circuit promotion is strongly backed by our Udaipur and Jaipur data -- those cities genuinely have the density and quality to justify the marketing spend. But Pushkar, which our Z-score analysis flags as undervalued with strong consumer engagement, gets almost zero campaign visibility despite being a FICCI-recognized wedding destination. This is the kind of misalignment that shows up only when you have structured data to compare against campaign investment.

A similar pattern likely exists in other states. Kerala promotes backwater weddings generically, but without venue-level price and capacity data, there is no way to know which specific micro-locations are underserved or overpriced. Our pipeline, extended to Kerala, would answer that question quantitatively.

\subsection{Overseas and NRI Market Report}

The NRI overseas wedding segment is estimated at USD 5 billion and growing. These clients have specific requirements: heritage aesthetics that photograph well for international social media audiences, high-capacity logistics (NRI weddings tend to be larger), transparent vendor contracting, and -- critically -- value for money. Many NRI families are price-sensitive despite high incomes because they compare Indian venue costs against international alternatives in Bali, Dubai, or Italy.

WedVenueIQ surfaced what we call ``Heritage-Value'' clusters -- venues in Pushkar and the outskirts of Jodhpur that offer palace-grade architectural aesthetics, 4.5+ ratings, and capacities above 300 guests, but at 30 to 40 percent lower prices than Udaipur's Tier-1 corridor. This is the exact intelligence that destination wedding planners working with NRI clients need. Currently, they operate on reputation and word-of-mouth; our dataset gives them data.

\subsection{Comparison with Existing Solutions}

\begin{table}[htbp]
\caption{Capability Matrix: Platforms vs. WedVenueIQ}
\begin{center}
\begin{tabular}{|l|c|c|}
\hline
\textbf{Capability} & \textbf{Platforms} & \textbf{WedVenueIQ} \\
\hline
Cross-platform aggregation & No & Yes \\
\hline
Numerical price (INR) & No & Yes \\
\hline
Deduplicated dataset & No & Yes \\
\hline
Luxury corridor detection & No & Yes \\
\hline
Undervalued city flagging & No & Yes \\
\hline
Micro-location scoring & No & Yes \\
\hline
Machine-readable export & No & Yes \\
\hline
Reproducible methodology & No & Yes \\
\hline
NRI market intelligence & No & Yes \\
\hline
\end{tabular}
\label{tab:comparison}
\end{center}
\end{table}

\section{Future Work}

\subsection{Machine Learning for Classification and Imputation}

The keyword-based venue type classifier leaves 16.8 percent of records unclassified. But here is the thing -- the 83.2 percent that it did classify correctly now form a labeled training set. A TF-IDF logistic regression trained on that set should push coverage above 95 percent with about 92 percent accuracy on held-out data. If we want to go further, fine-tuning a BERT model on Indian hospitality text (venue names, descriptions, type labels) could plausibly exceed 97 percent accuracy, though inference cost goes up.

On the imputation side, city-level median is crude. A RandomForestRegressor taking venue type, capacity, rating, review count, and city as features would produce venue-specific price estimates that capture interactions between those variables. We plan to measure RMSE improvement on a 20 percent validation set with known prices.

\subsection{Seasonal Pricing and Muhurat Integration}

The pipeline was designed for repeated runs. Execute it weekly, timestamp each scrape, and you get longitudinal price data. Our working hypothesis: venue prices in Rajasthan inflate by 15 to 25 percent during the November-to-February peak season when Muhurat dates cluster heavily. Nobody has tested this quantitatively. Integrating the Hindu Panchanga calendar to compute a weekly Muhurat density score would let us model auspicious-date effects on pricing explicitly. That intersection -- Indian cultural calendar data meets hospitality economics -- could be a genuinely novel contribution.

\subsection{Visual Feature Extraction}

Venue photos contain information that text listings miss entirely. A pool. A garden. Heritage sandstone architecture. A modern glass banquet interior. CLIP-based zero-shot classification can extract these features from scraped images as binary flags -- has\_pool, has\_garden, heritage\_architecture. Combined with text-based features, this multimodal approach should improve both type classification accuracy and price prediction. We have not implemented this yet but the image URLs are already scraped and stored.

\subsection{National Expansion}

The cities.json config file already lists 25 cities across 10 states: Rajasthan, Goa, Kerala, Maharashtra, Karnataka, Uttarakhand, Tamil Nadu, Himachal Pradesh, Delhi NCR, West Bengal. Running the pipeline across all of them should yield between 15,000 and 25,000 unique venue records. That would make this the largest structured Indian wedding venue dataset in existence by a wide margin.

Geocoding through the Nominatim API would add latitude-longitude to each record, unlocking spatial analytics: DBSCAN clustering of venue distributions, drive-time analysis from major airports, distance-to-heritage-site scoring for tourism positioning.

\subsection{API and Dashboard Deployment}

The planned deployment target is a FastAPI REST endpoint:

\begin{verbatim}
GET /api/v1/venues?state=Rajasthan
  &type=palace&max_price=15000000
  &min_capacity=200
\end{verbatim}

This returns a JSON array of matching venues with all normalized fields. A companion Streamlit web app would provide interactive map visualization, WVVI-ranked results, and city comparison dashboards. Together these would give the ``Wed in India'' campaign the kind of data backbone that tourism.gov.in currently lacks.

\section{Conclusion}

WedVenueIQ demonstrates that structured intelligence can be pulled from the fragmented Indian wedding market -- a sector the Prime Minister himself flagged as nationally important through the ``Wed in India'' call in Mann Ki Baat \cite{b22}. The market is worth 3.75 lakh crore INR \cite{b1}. The Government and FICCI have invested in campaigns, expos, influencer outreach \cite{b23} \cite{b24}. What was missing is data. Structured, normalized, deduplicated, queryable data.

The Rajasthan pilot proves the pipeline works. 5,942 raw records scraped from three heterogeneous sources. 16.3 percent cross-source duplicates removed through city-gated fuzzy matching. Indian-specific pricing strings normalized at 94 percent accuracy via a six-pass regex cascade. 2,036 clean, analytics-ready records produced in about 10 minutes total. These are not theoretical numbers; the dataset exists and the code runs.

The analytical layer on top -- WVVI for venue ranking, Z-score classification for city-level market typing, MLS for micro-location scoring -- shows that the dataset enables insights that were previously out of reach. Luxury corridors in Udaipur, undervalued destinations in Pushkar and Kumbhalgarh, hidden high-value micro-locations in Jaipur and Jodhpur. The NRI overseas analysis connects directly to the Government's cultural soft power agenda \cite{b21}: Indian weddings are not just ceremonies, they are a globally visible expression of heritage, and the data infrastructure to optimize their economic impact should exist.

What we have built is not a consumer app or a government portal. It is the data layer underneath both -- the structured foundation that campaigns, planners, researchers, and AI systems can build upon. The methodology is modular, reproducible, configuration-driven, and directly transferable to adjacent verticals: corporate events, banquet bookings, luxury domestic travel.

\section*{Acknowledgment}

The authors thank the faculty of the Department of Computer Science for guidance through this project. Data was collected from publicly accessible web listings in accordance with platform terms of service.

\begin{thebibliography}{00}
\bibitem{b1} KPMG India, ``Indian Wedding Industry Report 2023,'' KPMG LLP, New Delhi, 2023.
\bibitem{b2} A. Doan, A. Halevy, and Z. Ives, \textit{Principles of Data Integration}. Morgan Kaufmann, 2012.
\bibitem{b3} Y. Zhai and B. Liu, ``Web data extraction based on partial tree alignment,'' in \textit{Proc. WWW}, 2005, pp. 76--85.
\bibitem{b4} Selenium Project, ``Selenium WebDriver Documentation,'' 2024.
\bibitem{b5} Microsoft, ``Playwright for Python Documentation,'' 2024.
\bibitem{b6} F. Nwanganga and M. Chapple, \textit{Practical Machine Learning in Python}. Apress, 2020.
\bibitem{b7} C. Bizer, M. Hepp, and A. Harth, ``Schema.org as a foundation for web-scale data integration,'' \textit{J. Web Semantics}, vol. 65, 2021.
\bibitem{b8} Preliminary audit of Indian wedding platform structured data coverage, October 2024 (unpublished).
\bibitem{b9} P. Christen, \textit{Data Matching}. Springer, 2012.
\bibitem{b10} I. P. Fellegi and A. B. Sunter, ``A theory for record linkage,'' \textit{J. Amer. Stat. Assoc.}, vol. 64, no. 328, pp. 1183--1210, 1969.
\bibitem{b11} V. I. Levenshtein, ``Binary codes capable of correcting deletions, insertions, and reversals,'' \textit{Soviet Physics Doklady}, vol. 10, no. 8, pp. 707--710, 1966.
\bibitem{b12} W. Cohen, P. Ravikumar, and S. Fienberg, ``A comparison of string metrics for matching names and records,'' in \textit{Proc. KDD Workshop}, 2003.
\bibitem{b13} R. C. Steorts, R. Hall, and S. E. Fienberg, ``A Bayesian approach to graphical record linkage,'' \textit{J. Amer. Stat. Assoc.}, vol. 111, no. 516, pp. 1660--1672, 2016.
\bibitem{b14} G. Papadakis et al., ``Meta-blocking: Taking entity resolution to the next level,'' \textit{IEEE Trans. Knowl. Data Eng.}, vol. 26, no. 8, pp. 1946--1960, 2014.
\bibitem{b15} S. Mudgal et al., ``Deep learning for entity matching,'' in \textit{Proc. ACM SIGMOD}, 2018, pp. 19--34.
\bibitem{b16} M. J. Cafarella et al., ``WebTables: Exploring the power of tables on the web,'' \textit{Proc. VLDB Endowment}, vol. 1, no. 1, pp. 538--549, 2008.
\bibitem{b17} STR Global, ``Hotel Performance Data and Benchmarking,'' 2024.
\bibitem{b18} T. D. Andersson and D. Getz, ``Festival and event management in Nordic countries,'' \textit{Tourism Management}, vol. 30, no. 6, pp. 872--880, 2009.
\bibitem{b19} M. Ester et al., ``A density-based algorithm for discovering clusters,'' in \textit{Proc. KDD}, 1996, pp. 226--231.
\bibitem{b20} K. L. Xie, Z. Zhang, and Z. Zhang, ``Business value of online consumer reviews,'' \textit{Int. J. Hospitality Management}, vol. 43, pp. 1--12, 2014.
\bibitem{b21} S. Banerjee, \textit{Make Me Indian: Wedding Mandap Design and Cultural Identity}. Routledge India, 2021.
\bibitem{b22} Prime Minister's Office, ``Mann Ki Baat, Episode 107,'' November 26, 2023. Available: pmindia.gov.in
\bibitem{b23} FICCI and Ministry of Tourism, ``Wed in India Expo,'' Great India Travel Bazaar, Jaipur, May 2024.
\bibitem{b24} Ministry of Tourism, Government of India, ``India Says I Do -- Wedding Destination Campaign,'' tourism.gov.in, 2024.
\end{thebibliography}

\vspace{12pt}
\color{red}
IEEE conference templates contain guidance text for composing and formatting conference papers. Please ensure that all template text is removed from your conference paper prior to submission to the conference. Failure to remove the template text from your paper may result in your paper not being published.

\end{document}

